% Compile with LuaLaTeX (the appendix shows Lean source with Unicode characters):
%   lualatex main && lualatex main
\documentclass[11pt]{article}
\usepackage{fontspec}
\directlua{luaotfload.add_fallback("leanfallback", {"Apple Symbols:mode=harf;", "STIX Two Math:mode=harf;"})}
\setmonofont{Menlo}[Scale=0.82, RawFeature={fallback=leanfallback}]
\usepackage[margin=1in]{geometry}
\usepackage{amsmath,amssymb,amsthm,mathtools}
\usepackage{microtype}
\usepackage{fancyvrb}
\usepackage[colorlinks=true,linkcolor=blue,citecolor=blue,urlcolor=blue]{hyperref}
\urlstyle{same}

\newtheorem{theorem}{Theorem}
\newtheorem{lemma}[theorem]{Lemma}
\newtheorem{corollary}[theorem]{Corollary}
\theoremstyle{remark}
\newtheorem{remark}[theorem]{Remark}

\newcommand{\Dcost}{D}
\DeclareMathOperator{\diag}{diag}

\title{A sharp spectral-deficit stability inequality\\ for triangle-free graphs}
\author{Ryosuke Shiroshita\\[3pt]
\small\href{mailto:Ryosuke.Shiroshita@exactory.ai}{Ryosuke.Shiroshita@exactory.ai}}
\date{25 October 2026}
\hypersetup{pdftitle={A sharp spectral-deficit stability inequality for triangle-free graphs},pdfauthor={Ryosuke Shiroshita}}

\begin{document}
\maketitle

\begin{abstract}
Let $G$ be a finite simple triangle-free graph with $m$ edges and adjacency spectral radius
$\lambda$, and let $D(G)$ be the least number of edge insertions and deletions that turn $G$ into a
complete bipartite graph together with isolated vertices. Nosal's inequality states that
$\lambda^2\le m$. We prove that $D(G)\le 3(m-\lambda^2)$. The five-cycle attains equality, so the
constant $3$ is best possible. For bipartite graphs this inequality is a lemma of Li, Liu and
Ning; the theorem removes the bipartiteness assumption. The proof averages the costs of the
targets $K_{N(v),N(u)}$ over the edges $uv$ of $G$ with the weights $x_u^2/d(u)+x_v^2/d(v)$, where
$x$ is a unit eigenvector for $\lambda$, and bounds this average by two Cauchy--Schwarz
inequalities and a counting lemma that uses triangle-freeness. The theorem and the sharpness
example are formalized in Lean~4 with Mathlib; the proof is accepted by Lean's kernel and by five
further proof checkers.
\end{abstract}

\section{Introduction}

For a finite simple graph $G$ we write $e(G)$ for its number of edges and $\lambda(G)$ for the
spectral radius of its adjacency matrix ($\lambda(G)=0$ when $G$ has no vertices). Nosal
\cite{Nosal} proved that every triangle-free graph satisfies $\lambda(G)^2\le e(G)$; see also
Nikiforov \cite{Nikiforov}. Complete bipartite graphs attain equality,
since $\lambda(K_{s,t})^2=st=e(K_{s,t})$, and adding isolated vertices changes neither side. We
measure the distance from this extremal situation by the \emph{spectral deficit}
\[
  \xi(G)=e(G)-\lambda(G)^2\ \ge\ 0 ,
\]
and the distance from complete bipartite graphs by
\[
  D(G)=\min_{\substack{A,B\subseteq V(G)\\ A\cap B=\varnothing}}\bigl|E(G)\,\triangle\,E(K_{A,B})\bigr| ,
\]
where $K_{A,B}$ is the graph on $V(G)$ whose edges are the pairs with one end in $A$ and the
other in $B$, and empty parts are allowed. Thus $D(G)$ is the least number of edge insertions and
deletions that turn $G$ into a complete bipartite graph together with isolated vertices (the
graph without edges included).

Li, Liu and Zhang \cite[Theorem~4.1]{LLZ} proved a qualitative edge-spectral stability theorem for
triangle-free graphs. For bipartite graphs, Li, Liu and Ning \cite[Lemma~4.3]{LLN} proved the
linear bound $D(H)\le 3\,\xi(H)$; their target is $K_{N(v),N(u)}$ for a suitable edge $uv$, found
by averaging over all edges. The question whether the same coefficient holds for every
triangle-free graph was posed in \cite{Problem}, together with the bound $D(G)\le 15\,\xi(G)$. We
answer it.

\begin{theorem}\label{thm:main}
For every finite simple triangle-free graph $G$,
\[
  D(G)\ \le\ 3\bigl(e(G)-\lambda(G)^2\bigr).
\]
\end{theorem}

The constant cannot be lowered. Label the five-cycle $C_5$ cyclically by $1,\dots,5$. Then
$e(C_5)=5$ and $\lambda(C_5)=2$, so $\xi(C_5)=1$, and $D(C_5)=3$: the target $K_{\{1\},\{2,5\}}$
costs the three edges $23$, $34$, $45$, and no target costs less (Section~\ref{sec:sharp}). The
balanced blow-ups $C_5[t]$ attain equality as well (Remark~\ref{rem:equality}).

The proof gives slightly more: when $G$ has an edge, there is an edge $uv$ whose target
$K_{N(v),N(u)}$ costs at most $3\,\xi(G)$, as in the bipartite lemma of \cite{LLN}. Since
$D(G)=0$ only for complete bipartite graphs with isolated vertices, Theorem~\ref{thm:main} also
contains the equality case of Nosal's inequality: $\xi(G)=0$ forces $G$ to be such a graph.

\paragraph{Idea of the proof.} For an edge $uv$ of a triangle-free graph the neighbourhoods
$N(u)$ and $N(v)$ are disjoint, so $K_{N(v),N(u)}$ is a target. Li, Liu and Ning average its cost
uniformly over the edges, which works for bipartite graphs. For triangle-free graphs the uniform
average can exceed $3\xi$ (Remark~\ref{rem:weights}). We choose a vertex $u$ with probability
$x_u^2$, where $x$ is a unit eigenvector for $\lambda$, and then a uniformly random neighbour $v$
of $u$. Two Cauchy--Schwarz inequalities for $\|Ax\|^2=\lambda^2$ and one counting lemma
(Lemma~\ref{lem:count}), the only place where triangle-freeness is used beyond the disjointness of
$N(u)$ and $N(v)$, show that the expected cost is at most $3\xi$.

\section{Notation and the edge targets}

Let $G$ be a triangle-free graph on the vertex set $V$ with $m$ edges and adjacency matrix $A$.
We write $d(u)$ for the degree of $u$, $d=A\mathbf 1$ for the degree vector, $N(u)$ for the
neighbourhood of $u$ and $N[u]=N(u)\cup\{u\}$. For $S\subseteq V$ and $z\in V$ let
$d_S(z)=|N(z)\cap S|$. For a vertex $u$ put
\[
  W_u=V\setminus N[u],\qquad Z_u=e\bigl(G[W_u]\bigr).
\]
The codegree $c(u,z)=|N(u)\cap N(z)|$ equals $(A^2)_{uz}$, and $c(u,u)=d(u)$. Hence
$(A^2d)_u=\sum_z c(u,z)\,d(z)$ and $(A^4)_{uu}=\sum_z c(u,z)^2$. Triangle-freeness gives two
facts that we use repeatedly: $N(u)$ is an independent set for every $u$, and $c(u,v)=0$ for every
edge $uv$.

Because $N(u)$ is independent, every edge that meets $N(u)$ meets it in exactly one end. An edge
misses $N(u)$ exactly when both its ends lie in $W_u$, since $u$ has no neighbour in $W_u$.
Therefore
\begin{equation}\label{eq:Ad}
  (Ad)_u=\sum_{v\in N(u)}d(v)=m-Z_u .
\end{equation}

For an edge $uv$ define the cost of the target $K_{N(v),N(u)}$,
\[
  \Dcost_{uv}=\bigl|E(G)\,\triangle\,E(K_{N(v),N(u)})\bigr| .
\]
As $N(u)\cap N(v)=\varnothing$, this target is admissible, and $D(G)\le \Dcost_{uv}$.

\begin{lemma}\label{lem:cost}
For every edge $uv$, $\Dcost_{uv}=m+d(u)d(v)-2(A^3)_{uv}$. For every vertex $u$ with $d(u)\ge1$,
\begin{equation}\label{eq:mean}
  \overline\Dcost_u:=\frac{1}{d(u)}\sum_{v\in N(u)}\Dcost_{uv}
  \;=\;2m-Z_u-\frac{2(A^4)_{uu}}{d(u)} .
\end{equation}
\end{lemma}

\begin{proof}
The target has $d(u)d(v)$ edges. An edge of $G$ belongs to it exactly when it has one end in
$N(u)$ and the other in $N(v)$. Since these sets are disjoint, each such edge is counted exactly
once in $\sum_{a\in N(u)}\sum_{b\in N(v)}A_{ab}=(A^3)_{uv}$. Hence
$\Dcost_{uv}=m+d(u)d(v)-2(A^3)_{uv}$. Summing over $v\in N(u)$, and using
$\sum_{v}A_{uv}(A^3)_{vu}=(A^4)_{uu}$ and \eqref{eq:Ad},
\[
  \sum_{v\in N(u)}\Dcost_{uv}=d(u)\,m+d(u)\,(Ad)_u-2(A^4)_{uu}
  =d(u)\,(2m-Z_u)-2(A^4)_{uu} .\qedhere
\]
\end{proof}

\section{A counting lemma}

\begin{lemma}\label{lem:count}
For every vertex $u$ with $d(u)\ge1$,
\[
  (A^2d)_u-(A^4)_{uu}
  =\sum_{zz'\in E(G[W_u])}\bigl(d_{N(u)}(z)+d_{N(u)}(z')\bigr)
  \ \le\ d(u)\,Z_u .
\]
\end{lemma}

\begin{proof}
We have $(A^2d)_u-(A^4)_{uu}=\sum_z c(u,z)\bigl(d(z)-c(u,z)\bigr)$. The term $z=u$ vanishes
because $c(u,u)=d(u)$, and the terms $z\in N(u)$ vanish because $c(u,z)=0$ for an edge $uz$. For
$z\in W_u$ the vertex $z$ is not adjacent to $u$, so $N(z)\subseteq N(u)\cup W_u$; hence
$c(u,z)=d_{N(u)}(z)$ and $d(z)-c(u,z)=d_{W_u}(z)$. Therefore
\[
  (A^2d)_u-(A^4)_{uu}=\sum_{z\in W_u}d_{N(u)}(z)\,d_{W_u}(z)
  =\sum_{zz'\in E(G[W_u])}\bigl(d_{N(u)}(z)+d_{N(u)}(z')\bigr),
\]
where the second sum counts each edge $zz'$ of $G[W_u]$ once from each end. For such an edge,
triangle-freeness gives $N(z)\cap N(z')=\varnothing$, so
$d_{N(u)}(z)+d_{N(u)}(z')=|(N(z)\cup N(z'))\cap N(u)|\le d(u)$. Summing over the $Z_u$ edges of
$G[W_u]$ gives the bound.
\end{proof}

The slack of the lemma is $\sum_{zz'\in E(G[W_u])}|N(u)\setminus(N(z)\cup N(z'))|$.

\section{Two Cauchy--Schwarz bounds}

\begin{lemma}\label{lem:cs}
Let $G$ be a graph without isolated vertices and let $y\in\mathbb R^V$. Then
\begin{align}
  y^{\mathsf T}A^2y&\le\sum_u y_u^2\,(Ad)_u, \label{eq:cs1}\\
  y^{\mathsf T}A^2y&\le\sum_u y_u^2\,\frac{(A^2d)_u}{d(u)} . \label{eq:cs2}
\end{align}
\end{lemma}

\begin{proof}
We have
\[
  y^{\mathsf T}A^2y=\|Ay\|^2=\sum_v\Bigl(\sum_{u\in N(v)}y_u\Bigr)^2 .
\]
For \eqref{eq:cs1}, bound $\bigl(\sum_{u\in N(v)}y_u\bigr)^2\le d(v)\sum_{u\in N(v)}y_u^2$ and
exchange the order of summation:
\[
  \sum_v d(v)\sum_{u\in N(v)}y_u^2=\sum_u y_u^2\sum_{v\in N(u)}d(v)=\sum_u y_u^2\,(Ad)_u .
\]
For \eqref{eq:cs2}, use the weights $d(u)$:
\[
  \Bigl(\sum_{u\in N(v)}y_u\Bigr)^2
  \le\Bigl(\sum_{u\in N(v)}d(u)\Bigr)\Bigl(\sum_{u\in N(v)}\frac{y_u^2}{d(u)}\Bigr)
  =(Ad)_v\sum_{u\in N(v)}\frac{y_u^2}{d(u)},
\]
and exchange the order of summation:
$\sum_v(Ad)_v\sum_{u\in N(v)}y_u^2/d(u)=\sum_u\bigl(y_u^2/d(u)\bigr)\sum_{v\in N(u)}(Ad)_v$, and
$\sum_{v\in N(u)}(Ad)_v=(A^2d)_u$.
\end{proof}

\section{Proof of Theorem~\ref{thm:main}}\label{sec:proof}

If $m=0$, then $D(G)=0$ (take $A=B=\varnothing$) and $\lambda=0$, so the statement holds. Let
$m\ge1$. Deleting the isolated vertices changes neither $m$ nor $\lambda$. It does not change
$D(G)$ either: an isolated vertex can be removed from either part of a target without increasing
its cost, and every target of the smaller graph is a target of $G$. So we may assume that every
vertex has degree at least one.

The adjacency matrix is symmetric and nonnegative, so its spectral radius $\lambda$ is its
largest eigenvalue. Let $x$ be a real unit eigenvector for $\lambda$; any such vector works. Then
$x^{\mathsf T}A^2x=\|Ax\|^2=\lambda^2$ and $\sum_ux_u^2=1$. Adding \eqref{eq:cs1} and twice
\eqref{eq:cs2} at $y=x$, and bounding $(A^2d)_u\le(A^4)_{uu}+d(u)Z_u$ by Lemma~\ref{lem:count},
\[
  3\lambda^2\le\sum_ux_u^2\Bigl[(Ad)_u+\frac{2(A^4)_{uu}}{d(u)}+2Z_u\Bigr]
  =\sum_ux_u^2\Bigl[m+Z_u+\frac{2(A^4)_{uu}}{d(u)}\Bigr],
\]
where the equality is \eqref{eq:Ad}. By \eqref{eq:mean},
\[
  \sum_ux_u^2\,\overline\Dcost_u
  =\sum_ux_u^2\Bigl[2m-Z_u-\frac{2(A^4)_{uu}}{d(u)}\Bigr]
  =3m-\sum_ux_u^2\Bigl[m+Z_u+\frac{2(A^4)_{uu}}{d(u)}\Bigr]
  \le3m-3\lambda^2 .
\]
The left side is a convex combination of the numbers $\overline\Dcost_u$, and each
$\overline\Dcost_u$ is the mean of $\Dcost_{uv}$ over the nonempty set $N(u)$. Hence some vertex
$u$ has $\overline\Dcost_u\le3\xi$, and some neighbour $v$ of $u$ has
$\Dcost_{uv}\le\overline\Dcost_u\le3\xi$. Therefore $D(G)\le\Dcost_{uv}\le3\xi(G)$. \qed

\medskip
In terms of edges, the inequality just proved says that
\[
  \sum_{uv\in E(G)}\Bigl(\frac{x_u^2}{d(u)}+\frac{x_v^2}{d(v)}\Bigr)\Dcost_{uv}\ \le\ 3\,\xi(G),
\]
where the weights are nonnegative and sum to $1$.

\section{Sharpness}\label{sec:sharp}

For the five-cycle, $m=5$ and the all-ones vector is a positive eigenvector with eigenvalue $2$,
so $\lambda(C_5)=2$ and $\xi(C_5)=1$. For disjoint $A,B\subseteq V(C_5)$ with $a=|A|$ and
$b=|B|$, the cost of $K_{A,B}$ is $5+ab-2e(A,B)$, where $e(A,B)$ is the number of edges between
$A$ and $B$. We show that $2e(A,B)-ab\le2$, which gives $D(C_5)\ge3$. We may assume $a\le b$, so
$a\le2$ since $a+b\le5$. If $a=0$ the claim is clear. If $a=1$, then $e(A,B)\le\min(2,b)$ and
$2\min(2,b)-b\le2$. If $a=2$, then $e(A,B)\le4$ because every vertex has degree two, and
$e(A,B)\le3$ when $b=2$ because $C_5$ has no four-cycle; so $2e(A,B)-ab\le\max(6-4,\,8-6)=2$.
The target $K_{\{1\},\{2,5\}}$ costs $3$. Hence $D(C_5)=3=3\,\xi(C_5)$, and no coefficient below
$3$ bounds $D$ by a multiple of $\xi$ on all triangle-free graphs.

\section{Remarks}

\begin{remark}[Other weightings]\label{rem:weights}
The uniform average of $\Dcost_{uv}$ over the edges, which proves the bipartite lemma of
\cite{LLN}, can exceed $3\xi$ for triangle-free graphs. A search over all labelled triangle-free
graphs on at most seven vertices, with eigenvalues computed to $50$ digits, finds no example on
five or fewer vertices. On six vertices the largest excess occurs for the blow-up
$C_5[1,1,2,1,1]$ of the five-cycle in which one vertex is replaced by two nonadjacent copies.
There $m=7$ and the seven edge costs are $5,5,5,3,3,3,3$, so the uniform average is $27/7$. The
characteristic polynomial is $z(z^2+z-1)(z^3-z^2-5z+4)$, and $\lambda$ is the largest root $r$ of
the cubic factor. The excess is $27/7-3(7-r^2)=3r^2-120/7\approx0.0133$. It is positive because
the cubic is increasing on $[2,\infty)$ and takes the value $-12/7+10\sqrt{70}/49<0$ at
$\sqrt{40/7}$, so $r^2>40/7$.

Weights proportional to $x_ux_v$ also fail. In the double star $S(8,7)$, two adjacent vertices
$a$ and $b$ with $8$ and $7$ pendant neighbours, the edge costs are $7$ on the pendant edges at
$a$, $8$ on the pendant edges at $b$, and $56$ on the edge $ab$. Here $\lambda^2=8+2\sqrt2$ and
$3\xi=24-6\sqrt2$, while $\tfrac{2}{\lambda}\sum_{uv\in E}x_ux_v\Dcost_{uv}=15/2+6\sqrt2$, larger by
$12\sqrt2-33/2\approx0.47$. The bound $D\le3\xi$ itself holds there with room, since a pendant edge
at $a$ gives a target of cost $7$.
\end{remark}

\begin{remark}[A matrix form]
Lemmas~\ref{lem:count} and~\ref{lem:cs} show that for every triangle-free graph without isolated
vertices the matrix
\[
  Q=\diag\Bigl(m+Z_u+\frac{2(A^4)_{uu}}{d(u)}\Bigr)_{u\in V}-3A^2
\]
is positive semidefinite; the key inequality of Section~\ref{sec:proof} is $x^{\mathsf T}Qx\ge0$.
\end{remark}

\begin{remark}[Equality]\label{rem:equality}
The three inequalities used in Section~\ref{sec:proof} are all equalities for complete bipartite
graphs and for the balanced blow-ups $C_5[t]$, in which every vertex of $C_5$ is replaced by an
independent set of size $t$ and every edge by a complete bipartite graph $K_{t,t}$. On $K_{s,t}$
the eigenvector is constant on each part and $G[W_u]$ has no edges. On the regular graph $C_5[t]$
the eigenvector is constant, and for every $u$ each edge $zz'$ of $G[W_u]$ joins the two classes
at distance two from the class of $u$, whose neighbourhoods together contain $N(u)$. Hence the
weighted average equals $3\xi$ on these graphs.

Equality in Theorem~\ref{thm:main} itself also holds on every $C_5[t]$: $D(C_5[t])=3t^2=3\,\xi(C_5[t])$
for every $t\ge1$. This follows from $\xi(H[t])=t^2\xi(H)$ and $D(H[t])=t^2D(H)$, which hold for
every graph $H$ and every $t\ge1$, where $H[t]$ replaces each vertex of $H$ by an independent set
of size $t$ and each edge by a complete bipartite graph $K_{t,t}$. The first identity holds
because $e(H[t])=t^2e(H)$ and the adjacency matrix of $H[t]$ is the Kronecker product of that of
$H$ with the all-ones $t\times t$ matrix, so its nonzero eigenvalues are $t$ times those of $H$.
For the second, write the cost of a target as $e(H[t])-g(A,B)$ with the gain
$g(A,B)=2e(A,B)-|A||B|$. Fix $A$ and $B$ outside one class $X$, and let $a=|A\cap X|$ and
$b=|B\cap X|$, so $a+b\le t$. Since $X$ is independent, the gain has the form
$\gamma+\alpha a+\beta b-ab$, which is at most $\gamma+t\max(\alpha,\beta,0)$; this value is attained
by putting all of $X$ into $A$, all of $X$ into $B$, or $X$ into neither part. Rounding the classes
one at a time gives an optimal target that is a union of whole classes, and the gain of such a
target is $t^2$ times the gain of the corresponding target of $H$.
\end{remark}

\section{Formal verification}

Theorem~\ref{thm:main} and the sharpness example are formalized in Lean~4 \cite{Lean4} with
Mathlib \cite{Mathlib} (Lean v4.35.0, Mathlib v4.35.0). The formal statement uses the following
definitions, reproduced in Appendix~\ref{app:lean}.
\begin{itemize}
  \item \texttt{editCost G S T} is $|E(G)\triangle E(K_{S,T})|$ on the vertex type, and
  \texttt{editDistance G} is its minimum over all pairs of disjoint finite sets $(S,T)$, empty
  sets allowed; this is $D(G)$.
  \item \texttt{spectralRadius G} is the largest absolute value of an eigenvalue of the real
  adjacency matrix ($0$ when there are no vertices). For an adjacency matrix this is the largest
  eigenvalue $\lambda(G)$.
  \item \texttt{G.CliqueFree 3} says that $G$ has no triangle. The vertex type is any finite type
  with decidable equality, so disconnected graphs and isolated vertices are included.
\end{itemize}
Two theorems state the results.
\begin{itemize}
  \item \texttt{TriangleFreeEdit.editDistance\_le\_three\_deficit}\\
  For every such $G$,
  $(\texttt{editDistance G}:\mathbb R)\le3\,(|E(G)|-\texttt{spectralRadius G}^2)$.
  \item \texttt{TriangleFreeSharpness.triangleFree\_constant\_sharp}\\
  For every real $c<3$ some finite triangle-free graph has $D(G)>c\,\xi(G)$.
\end{itemize}
For both, the command \texttt{\#print axioms} reports only \texttt{propext},
\texttt{Classical.choice} and \texttt{Quot.sound}.

The formal statement of the problem is a separate file, \texttt{Challenge.lean}, which contains
the definitions and the theorem with the proof replaced by \texttt{sorry}. The tool
\texttt{lake comparator} of Lean~v4.35.0 checks that the proved theorem has the same statement as
the challenge, uses only the permitted axioms, and is accepted by the kernel. With the option
\texttt{--paranoid} it also replays the proof in five further kernels (Lean's paranoid kernel,
\texttt{lean4lean}, \texttt{nanoda}, \texttt{con-leche} and \texttt{con-ron}); all six kernels
accepted it on 25 October 2026. That run used the option \texttt{--inadvisably-no-sandbox},
because the sandbox requires Linux; a checker who does not trust our files can repeat it with the
sandbox. The four external checkers are those bundled with the Lean v4.35.0 toolchain.

To reproduce the checks, run in the folder \texttt{lean/} of the repository (see Data
availability): \texttt{lake exe cache get}, \texttt{lake build}, \texttt{lake env lean Audit.lean}
and \texttt{lake comparator --paranoid}. The output of our run is the file
\texttt{lean/checks/comparator-paranoid.txt}.

\section*{How the result was obtained}

The problem was created by the math-conjecturer harness of exactory and posted on 22 October 2026
\cite{Problem}. The proof was found on 24 October 2026 by the math-solver harness of exactory, in
which the Anthropic language model Claude Fable~5.1 wrote the proof, Claude Opus~5.5 reviewed it
independently, and the OpenAI model gpt-6-astra reviewed it adversarially. Every identity and inequality of the
proof was also checked numerically, with exact integer arithmetic for the combinatorial steps, on
$2{,}918$ triangle-free graphs (the script \texttt{record/scripts/verify\_proof\_steps.py}, with
its output in \texttt{record/checks/}). The Lean proof was written by OpenAI Codex with gpt-6-astra and checked as
described above. No human mathematician has yet reviewed the proof.

\section*{Data availability}

The Lean sources, the formal statement, the configuration of \texttt{lake comparator}, the logs
of the checks and the scripts of the numerical checks are in the public repository
\url{https://github.com/exactory/exactory-math}, folder
\texttt{attempts/triangle-free-spectral-stability}. The version checked as described in this paper
is commit \texttt{7194c5784dba0cfe43dd5951f65e62f1b9c9456e}: the Lean files are in \texttt{lean/},
the log of \texttt{lake comparator} in \texttt{lean/checks/}, and the scripts of the numerical
checks of Remark~\ref{rem:weights} and of the proof steps, with their outputs, in \texttt{record/scripts/} and
\texttt{record/checks/}.

\begin{thebibliography}{9}
\bibitem{Lean4} L.~de~Moura and S.~Ullrich, The Lean~4 theorem prover and programming language,
  in: \emph{Automated Deduction -- CADE 28}, Lecture Notes in Computer Science 12699, Springer,
  2021, 625--635.
\bibitem{Problem} exactory, Sharp spectral-deficit stability for triangle-free graphs, open
  problem posted 22 October 2026,
  \url{https://www.exactory.ai/math/5aa390a0-bf63-4326-9b97-b50c46334fd8}.
\bibitem{LLN} Y.~Li, L.~Liu and B.~Ning, The spectral Erd\H{o}s book theorem: sharp bounds and
  stability, arXiv:2609.20225v2.
\bibitem{LLZ} Y.~Li, H.~Liu and S.~Zhang, An edge-spectral Erd\H{o}s--Stone--Simonovits theorem
  and its stability, arXiv:2508.15271v2.
\bibitem{Mathlib} The mathlib Community, The Lean mathematical library, in: \emph{Proceedings of
  the 9th ACM SIGPLAN International Conference on Certified Programs and Proofs (CPP 2020)}, 2020,
  367--381.
\bibitem{Nikiforov} V.~Nikiforov, On a theorem of Nosal, arXiv:2104.12171.
\bibitem{Nosal} E.~Nosal, Eigenvalues of graphs, Master's thesis, University of Calgary, 1970.
\end{thebibliography}

\appendix
\section{The formal statement}\label{app:lean}

The definitions and the theorem statement of \texttt{Challenge.lean}, without its header
comments. In \texttt{TriangleFreeEdit.lean} the same lines are followed by the proof.

\begin{Verbatim}[fontsize=\small]
namespace TriangleFreeEdit

variable {V : Type*} [Fintype V] [DecidableEq V]

def bicliqueEdges (S T : Finset V) : Finset (Sym2 V) :=
  (S ×ˢ T).image (fun p => s(p.1, p.2))

def editCost (G : SimpleGraph V) [DecidableRel G.Adj] (S T : Finset V) : ℕ :=
  (symmDiff G.edgeFinset (bicliqueEdges S T)).card

def disjointPairs (V : Type*) [Fintype V] [DecidableEq V] : Finset (Finset V × Finset V) :=
  (Finset.univ : Finset (Finset V × Finset V)).filter (fun p => Disjoint p.1 p.2)

lemma disjointPairs_nonempty : (disjointPairs V).Nonempty :=
  ⟨(∅, ∅), by simp [disjointPairs]⟩

lemma editCosts_nonempty (G : SimpleGraph V) [DecidableRel G.Adj] :
    ((disjointPairs V).image (fun p => editCost G p.1 p.2)).Nonempty :=
  Finset.Nonempty.image disjointPairs_nonempty _

noncomputable def editDistance (G : SimpleGraph V) [DecidableRel G.Adj] : ℕ :=
  ((disjointPairs V).image (fun p => editCost G p.1 p.2)).min' (by exact editCosts_nonempty G)

lemma adjMatrix_isHermitian (G : SimpleGraph V) [DecidableRel G.Adj] :
    (G.adjMatrix ℝ).IsHermitian := by
  exact Matrix.isHermitian_iff_isSymm.mpr (G.isSymm_adjMatrix)

noncomputable def spectralRadius (G : SimpleGraph V) [DecidableRel G.Adj] : ℝ :=
  if h : Nonempty V then
    (Finset.univ : Finset V).sup' Finset.univ_nonempty
      (fun i => |(adjMatrix_isHermitian G).eigenvalues i|)
  else 0

theorem editDistance_le_three_deficit (G : SimpleGraph V) [DecidableRel G.Adj]
    (hG : G.CliqueFree 3) :
    (editDistance G : ℝ) ≤ 3 * ((G.edgeFinset.card : ℝ) - (spectralRadius G) ^ 2) := by
  sorry

end TriangleFreeEdit
\end{Verbatim}

\end{document}
